\documentclass[11pt,a4paper]{article}

% ============================================================
%  QIMMA -- COURS : Produit scalaire et produit vectoriel
%  dans l'espace
%  2ème Bac Sciences Mathématiques (A et B)
%  Standard exhaustif : définitions formelles + démonstrations
%  Basé sur templates/qimma-template.tex (charte Qimma)
% ============================================================

\usepackage[french]{babel}
\usepackage[utf8]{inputenc}
\usepackage[T1]{fontenc}
\usepackage{lmodern}
\usepackage[a4paper,margin=2cm,top=2.3cm,bottom=2.6cm]{geometry}
\usepackage{amsmath,amssymb}
\usepackage{xcolor}
\usepackage{graphicx}
\usepackage{tikz}
\usetikzlibrary{shadows,calc}
\usepackage[most]{tcolorbox}
\usepackage{titlesec}
\usepackage{fancyhdr}
\usepackage{enumitem}
\usepackage{pifont}
\usepackage{array}
\usepackage{colortbl}
\usepackage{hyperref}

% ------------------- CHARTE QIMMA -------------------
\definecolor{qteal}{HTML}{0d9488}
\definecolor{qtealdark}{HTML}{134e4a}
\definecolor{qamber}{HTML}{fbbf24}
\definecolor{qambertext}{HTML}{92400e}
\definecolor{ink}{HTML}{1f2937}
\definecolor{inkgray}{HTML}{6b7280}
\definecolor{tealpale}{HTML}{e6f5f3}
\colorlet{colDef}{qteal}
\definecolor{colProp}{HTML}{0f766e}
\definecolor{colEx}{HTML}{b45309}
\definecolor{colRem}{HTML}{9d174d}
\color{ink}

\graphicspath{{../../../../assets/}}
\newcommand{\logofile}{../../../../assets/qimma-logo.pdf}

% ------------------- METADONNEES -------------------
\newcommand{\matiere}{Mathématiques}
\newcommand{\niveau}{2\up{ème} Bac -- Sciences Mathématiques}
\newcommand{\chapitretitre}{Produit scalaire et produit vectoriel dans l'espace}

% ------------------- EN-TETE / PIED DE PAGE -------------------
\pagestyle{fancy}
\fancyhf{}
\renewcommand{\headrulewidth}{0pt}
\renewcommand{\footrulewidth}{0.5pt}
\fancyfoot[L]{\raisebox{-0.15cm}{\includegraphics[height=0.35cm]{\logofile}}~\small\sffamily\color{inkgray}Qimma}
\fancyfoot[C]{\small\sffamily\color{inkgray}\matiere\ -- \niveau}
\fancyfoot[R]{\small\sffamily\color{qtealdark}\thepage}

% ------------------- TITRES DE SECTION -------------------
\titleformat{\section}
  {\normalfont\sffamily\Large\bfseries\color{qtealdark}}
  {\tikz[baseline=-3.2pt]{\node[circle,fill=qteal,text=white,inner sep=0pt,minimum size=8mm]{\sffamily\bfseries\thesection};}}
  {10pt}{}
  [\vspace{-2mm}{\color{qamber}\titlerule[1.6pt]}\vspace{2mm}]
\titlespacing*{\section}{0pt}{16pt}{10pt}

\titleformat{\subsection}
  {\normalfont\sffamily\large\bfseries\color{qtealdark}}
  {\color{qteal}\ding{212}}{6pt}{}
\titlespacing*{\subsection}{0pt}{12pt}{6pt}

% ------------------- BOITES PEDAGOGIQUES -------------------
\newtcolorbox{fiche}[3][]{
  enhanced, breakable,
  colback=white, colframe=#2,
  boxrule=1pt, arc=2mm,
  top=7mm, left=4mm, right=4mm, bottom=3mm,
  drop shadow={black!25, shadow xshift=1mm, shadow yshift=-1mm},
  attach boxed title to top left={xshift=4mm, yshift=-3mm, yshifttext=-1mm},
  boxed title style={colback=#2, arc=1.5mm, boxrule=0pt, left=2mm, right=2mm},
  coltitle=white, fonttitle=\sffamily\bfseries,
  title={#3}, #1
}
\newenvironment{definition}[1][Définition]
  {\begin{fiche}{colDef}{\ding{110}~~#1}}{\end{fiche}}
\newenvironment{propriete}[1][Propriété]
  {\begin{fiche}{colProp}{\ding{72}~~#1}}{\end{fiche}}
\newenvironment{exemple}[1][Exemple]
  {\begin{fiche}{colEx}{\ding{217}~~#1}}{\end{fiche}}
\newenvironment{remarque}[1][Remarque]
  {\begin{fiche}{colRem}{\ding{43}~~#1}}{\end{fiche}}

% Démonstration (hors boîte, style discret)
\newenvironment{demo}
  {\par\smallskip\noindent{\sffamily\bfseries\color{qtealdark}Démonstration.}\ \itshape}
  {\ \normalfont\hfill$\blacksquare$\par\medskip}

\hypersetup{colorlinks=true, linkcolor=qtealdark, urlcolor=qtealdark}

% ============================================================
\begin{document}

% ------------------- BANNIERE DE TITRE -------------------
\begin{tcolorbox}[
  enhanced, boxrule=0pt, arc=4mm,
  interior style={left color=qtealdark, right color=qteal},
  top=8mm, bottom=8mm, left=10mm, right=32mm,
  drop shadow={black!35, shadow xshift=1.5mm, shadow yshift=-1.5mm},
  before skip=0pt, after skip=9mm,
  overlay={
    \node[anchor=north west] at ([xshift=6mm,yshift=-6mm]frame.north west)
      {\includegraphics[height=1.25cm]{\logofile}};
    \node[anchor=north east, fill=qamber, text=qambertext, rounded corners=2pt,
          inner sep=4pt, font=\sffamily\bfseries\small] at ([xshift=-6mm,yshift=-6mm]frame.north east) {COURS};
  }
]
\hspace{1.6cm}\centering
{\sffamily\bfseries\color{white}\fontsize{20}{24}\selectfont \chapitretitre}\\[3mm]
{\sffamily\color{white!90}\large \matiere\ \textbullet\ \niveau}
\end{tcolorbox}

% ------------------- OBJECTIFS -------------------
\begin{tcolorbox}[
  enhanced, colback=tealpale, colframe=qteal, boxrule=0.8pt, arc=2mm,
  left=5mm, right=5mm, top=3mm, bottom=3mm,
  title={\sffamily\bfseries\color{qtealdark}\ding{51}~~Objectifs du chapitre},
  coltitle=qtealdark, colbacktitle=tealpale, boxrule=0.8pt,
  attach title to upper={\par\vspace{1mm}}
]
\begin{itemize}[leftmargin=6mm, label=\ding{212}, itemsep=1mm]
  \item Étendre le produit scalaire à l'espace, démontrer ses propriétés, et l'exprimer en coordonnées dans un repère orthonormé.
  \item Utiliser le produit scalaire pour caractériser l'orthogonalité, calculer normes, distances et angles.
  \item Définir le produit vectoriel, démontrer ses propriétés algébriques et son expression en coordonnées.
  \item Utiliser le produit vectoriel pour tester la colinéarité, calculer des aires de triangles/parallélogrammes et des volumes.
  \item Résoudre les exercices type examen national combinant les deux produits (vecteur normal à un plan, distance point-plan, etc.).
\end{itemize}
\end{tcolorbox}

\vspace{4mm}

% ------------------- SECTION 1 -------------------
\section{Produit scalaire dans l'espace}

\begin{definition}[Produit scalaire]
Pour $\vec u,\vec v$ vecteurs de l'espace non nuls, le \textbf{produit scalaire} est
\[
  \vec u\cdot\vec v = \|\vec u\|\,\|\vec v\|\,\cos(\vec u,\vec v).
\]
Si $\vec u=\vec0$ ou $\vec v=\vec0$ : $\vec u\cdot\vec v=0$.
\end{definition}

\begin{propriete}[Propriétés algébriques]
Pour tous vecteurs $\vec u,\vec v,\vec w$ de l'espace et $\lambda\in\mathbb{R}$ :
\[
  \vec u\cdot\vec v = \vec v\cdot\vec u, \qquad
  \vec u\cdot(\vec v+\vec w) = \vec u\cdot\vec v+\vec u\cdot\vec w, \qquad
  (\lambda\vec u)\cdot\vec v = \lambda(\vec u\cdot\vec v),
\]
\[
  \vec u\cdot\vec u = \|\vec u\|^2, \qquad
  \vec u\perp\vec v \iff \vec u\cdot\vec v=0.
\]
\end{propriete}

\begin{propriete}[Expression en repère orthonormé]
Dans un repère orthonormé de l'espace, si $\vec u(x,y,z)$ et $\vec v(x',y',z')$ :
\[
  \vec u\cdot\vec v = xx'+yy'+zz', \qquad \|\vec u\| = \sqrt{x^2+y^2+z^2}.
\]
\end{propriete}

\begin{demo}
Décomposons $\vec u=x\vec i+y\vec j+z\vec k$, $\vec v=x'\vec i+y'\vec j+z'\vec k$ où $(\vec i,\vec j,\vec k)$ est la base orthonormée du repère. Par bilinéarité du produit scalaire :
\[
  \vec u\cdot\vec v = \sum_{\text{termes}} (\text{coefficients})(\vec i\cdot\vec i,\ \vec i\cdot\vec j,\ldots).
\]
Comme $(\vec i,\vec j,\vec k)$ est orthonormée : $\vec i\cdot\vec i=\vec j\cdot\vec j=\vec k\cdot\vec k=1$ et $\vec i\cdot\vec j=\vec j\cdot\vec k=\vec i\cdot\vec k=0$. En développant, seuls les termes « carrés » (mêmes vecteurs de base) subsistent : $\vec u\cdot\vec v=xx'+yy'+zz'$.
\end{demo}

\begin{propriete}[Formule d'Al-Kashi vectorielle]
\[
  \|\vec u+\vec v\|^2 = \|\vec u\|^2+2\vec u\cdot\vec v+\|\vec v\|^2, \qquad
  \|\vec u-\vec v\|^2 = \|\vec u\|^2-2\vec u\cdot\vec v+\|\vec v\|^2.
\]
Conséquence géométrique (dans un triangle $ABC$) :
\[
  BC^2 = AB^2+AC^2-2\,AB\times AC\times\cos\widehat{A}.
\]
\end{propriete}

\begin{exemple}[Calcul de produit scalaire et angle]
Soient $A(1,0,2)$, $B(3,1,1)$, $C(2,2,0)$ dans un repère orthonormé. Calculons $\overrightarrow{AB}\cdot\overrightarrow{AC}$ et l'angle $\widehat{BAC}$.

\smallskip
$\overrightarrow{AB}(2,1,-1)$, $\overrightarrow{AC}(1,2,-2)$.
\[
  \overrightarrow{AB}\cdot\overrightarrow{AC} = 2\times1+1\times2+(-1)\times(-2) = 2+2+2 = 6.
\]
\[
  \|\overrightarrow{AB}\| = \sqrt{4+1+1}=\sqrt6, \qquad \|\overrightarrow{AC}\| = \sqrt{1+4+4}=3.
\]
\[
  \cos\widehat{BAC} = \frac{6}{\sqrt6\times3} = \frac{6}{3\sqrt6} = \frac{2}{\sqrt6} = \frac{2\sqrt6}{6} = \frac{\sqrt6}{3}.
\]
\end{exemple}

% ------------------- SECTION 2 -------------------
\section{Vecteur normal à un plan}

\begin{definition}[Vecteur normal]
Un vecteur $\vec n\ne\vec0$ est \textbf{normal} à un plan $\mathcal{P}$ lorsque $\vec n$ est orthogonal à tout vecteur directeur de $\mathcal{P}$ (de façon équivalente, orthogonal à deux vecteurs non colinéaires de $\mathcal{P}$).
\end{definition}

\begin{propriete}[Équation cartésienne d'un plan]
Le plan passant par $A(x_0,y_0,z_0)$ et de vecteur normal $\vec n(a,b,c)$ a pour équation cartésienne
\[
  a(x-x_0)+b(y-y_0)+c(z-z_0)=0, \qquad\text{soit}\qquad ax+by+cz+d=0.
\]
Réciproquement, tout plan d'équation $ax+by+cz+d=0$ (avec $(a,b,c)\ne(0,0,0)$) admet $\vec n(a,b,c)$ pour vecteur normal.
\end{propriete}

\begin{demo}
$M(x,y,z)\in\mathcal{P} \iff \overrightarrow{AM}$ est orthogonal à $\vec n$ (car $\overrightarrow{AM}$ est alors un vecteur directeur du plan, ou nul) $\iff \vec n\cdot\overrightarrow{AM}=0 \iff a(x-x_0)+b(y-y_0)+c(z-z_0)=0$.
\end{demo}

\begin{propriete}[Distance d'un point à un plan]
La distance du point $M_0(x_0,y_0,z_0)$ au plan $\mathcal{P}:ax+by+cz+d=0$ est
\[
  d(M_0,\mathcal{P}) = \frac{|ax_0+by_0+cz_0+d|}{\sqrt{a^2+b^2+c^2}}.
\]
\end{propriete}

\begin{exemple}[Équation de plan et distance]
Déterminons l'équation du plan $\mathcal{P}$ passant par $A(1,2,-1)$ de vecteur normal $\vec n(2,-1,3)$, puis la distance de $O(0,0,0)$ à $\mathcal{P}$.

\smallskip
Équation : $2(x-1)-1(y-2)+3(z+1)=0 \iff 2x-y+3z+3=0$.

\smallskip
Distance de $O$ : $d(O,\mathcal{P})=\dfrac{|2\times0-0+3\times0+3|}{\sqrt{4+1+9}}=\dfrac{3}{\sqrt{14}}=\dfrac{3\sqrt{14}}{14}$.
\end{exemple}

% ------------------- SECTION 3 -------------------
\section{Produit vectoriel}

\begin{definition}[Produit vectoriel]
Pour $\vec u,\vec v$ vecteurs de l'espace, le \textbf{produit vectoriel} $\vec u\wedge\vec v$ est le vecteur défini par :
\begin{itemize}[leftmargin=6mm, itemsep=0.5mm]
  \item si $\vec u,\vec v$ colinéaires (ou l'un est nul) : $\vec u\wedge\vec v=\vec0$ ;
  \item sinon : $\vec u\wedge\vec v$ est orthogonal à $\vec u$ et à $\vec v$, sa norme vaut $\|\vec u\|\,\|\vec v\|\,|\sin(\vec u,\vec v)|$, et le triplet $(\vec u,\vec v,\vec u\wedge\vec v)$ est direct.
\end{itemize}
\end{definition}

\begin{propriete}[Expression en coordonnées]
Dans une base orthonormée directe, si $\vec u(x,y,z)$, $\vec v(x',y',z')$ :
\[
  \vec u\wedge\vec v = \bigl(yz'-zy'\,;\ zx'-xz'\,;\ xy'-yx'\bigr).
\]
\end{propriete}

\begin{propriete}[Propriétés algébriques]
\[
  \vec u\wedge\vec v = -\vec v\wedge\vec u \quad\text{(antisymétrie)}, \qquad
  \vec u\wedge(\vec v+\vec w) = \vec u\wedge\vec v+\vec u\wedge\vec w,
\]
\[
  (\lambda\vec u)\wedge\vec v = \lambda(\vec u\wedge\vec v), \qquad
  \vec u\wedge\vec v = \vec0 \iff \vec u,\vec v \text{ colinéaires}.
\]
\end{propriete}

\begin{propriete}[Interprétation en aires]
$\|\vec u\wedge\vec v\|$ est l'\textbf{aire du parallélogramme} construit sur $\vec u$ et $\vec v$. L'aire du triangle correspondant est
\[
  \mathcal{A}_{ABC} = \frac12\left\|\overrightarrow{AB}\wedge\overrightarrow{AC}\right\|.
\]
\end{propriete}

\begin{exemple}[Calcul de produit vectoriel et aire]
Soient $\vec u(1,2,-1)$, $\vec v(3,0,1)$. Calculons $\vec u\wedge\vec v$ et l'aire du parallélogramme associé.
\[
  \vec u\wedge\vec v = \bigl(2\times1-(-1)\times0\,;\ (-1)\times3-1\times1\,;\ 1\times0-2\times3\bigr) = (2,-4,-6).
\]
\[
  \|\vec u\wedge\vec v\| = \sqrt{4+16+36} = \sqrt{56} = 2\sqrt{14}.
\]
L'aire du parallélogramme vaut $2\sqrt{14}$ (u.a.).
\end{exemple}

\begin{propriete}[Vecteur normal via le produit vectoriel]
Si $\vec u,\vec v$ sont deux vecteurs directeurs (non colinéaires) d'un plan $\mathcal{P}$, alors $\vec u\wedge\vec v$ est un vecteur \textbf{normal} à $\mathcal{P}$.
\end{propriete}

\begin{exemple}[Trouver un plan par trois points via le produit vectoriel]
Déterminons l'équation du plan $(ABC)$ où $A(1,0,0)$, $B(0,1,0)$, $C(0,0,1)$.

\smallskip
$\overrightarrow{AB}(-1,1,0)$, $\overrightarrow{AC}(-1,0,1)$.
\[
  \vec n = \overrightarrow{AB}\wedge\overrightarrow{AC} = (1\times1-0\times0\,;\ 0\times(-1)-(-1)\times1\,;\ (-1)\times0-1\times(-1)) = (1,1,1).
\]
Équation du plan (normal $(1,1,1)$, passant par $A(1,0,0)$) : $1(x-1)+1(y-0)+1(z-0)=0\iff x+y+z-1=0$.
\end{exemple}

\begin{propriete}[Produit mixte et volume]
Le \textbf{produit mixte} de trois vecteurs $\vec u,\vec v,\vec w$ est $\bigl(\vec u\wedge\vec v\bigr)\cdot\vec w$. Sa valeur absolue est le \textbf{volume du parallélépipède} construit sur ces trois vecteurs :
\[
  V_{\text{parallélépipède}} = \bigl|(\vec u\wedge\vec v)\cdot\vec w\bigr|, \qquad
  V_{\text{tétraèdre } ABCD} = \frac16\left|\bigl(\overrightarrow{AB}\wedge\overrightarrow{AC}\bigr)\cdot\overrightarrow{AD}\right|.
\]
Trois vecteurs sont \textbf{coplanaires} $\iff$ leur produit mixte est nul.
\end{propriete}

\begin{exemple}[Volume d'un tétraèdre]
Soient $A(0,0,0)$, $B(1,0,0)$, $C(0,1,0)$, $D(0,0,1)$. Calculons le volume du tétraèdre $ABCD$.

\smallskip
$\overrightarrow{AB}(1,0,0)$, $\overrightarrow{AC}(0,1,0)$, $\overrightarrow{AD}(0,0,1)$.
\[
  \overrightarrow{AB}\wedge\overrightarrow{AC} = (0\times0-0\times1\,;\ 0\times0-1\times0\,;\ 1\times1-0\times0) = (0,0,1).
\]
\[
  (\overrightarrow{AB}\wedge\overrightarrow{AC})\cdot\overrightarrow{AD} = 0\times0+0\times0+1\times1 = 1.
\]
\[
  V = \frac16|1| = \frac16.
\]
\end{exemple}

% ------------------- SECTION 4 -------------------
\section{Exercices corrigés -- niveau examen national SM}

\begin{exemple}[Exercice 1 -- orthogonalité et équation de plan]
Soient $A(1,1,1)$, $B(2,0,3)$. Déterminer l'équation du plan médiateur de $[AB]$ (plan orthogonal à $(AB)$ passant par son milieu).

\smallskip
\textbf{Corrigé.} $\overrightarrow{AB}(1,-1,2)$ est un vecteur normal du plan médiateur. Milieu $I\left(\dfrac32,\dfrac12,2\right)$.

\smallskip
Équation : $1\left(x-\dfrac32\right)-1\left(y-\dfrac12\right)+2(z-2)=0 \iff x-y+2z-\dfrac{3}{2}+\dfrac12-4=0 \iff x-y+2z-5=0$.
\end{exemple}

\begin{exemple}[Exercice 2 -- calcul d'aire par produit vectoriel]
Soient $A(1,0,0)$, $B(0,2,0)$, $C(0,0,3)$. Calculer l'aire du triangle $ABC$.

\smallskip
\textbf{Corrigé.} $\overrightarrow{AB}(-1,2,0)$, $\overrightarrow{AC}(-1,0,3)$.
\[
  \overrightarrow{AB}\wedge\overrightarrow{AC} = (2\times3-0\times0\,;\ 0\times(-1)-(-1)\times3\,;\ (-1)\times0-2\times(-1)) = (6,3,2).
\]
\[
  \left\|\overrightarrow{AB}\wedge\overrightarrow{AC}\right\| = \sqrt{36+9+4} = \sqrt{49} = 7.
\]
\[
  \mathcal{A}_{ABC} = \frac72.
\]
\end{exemple}

\begin{exemple}[Exercice 3 -- distance et projeté orthogonal]
Soit $\mathcal{P}:x+2y-2z+1=0$ et $M(1,2,3)$. Calculer $d(M,\mathcal{P})$, puis déterminer les coordonnées du projeté orthogonal $H$ de $M$ sur $\mathcal{P}$.

\smallskip
\textbf{Corrigé.} $d(M,\mathcal{P}) = \dfrac{|1+4-6+1|}{\sqrt{1+4+4}} = \dfrac{0}{3} = 0$ : le point $M$ appartient au plan $\mathcal{P}$ (vérification directe : $1+4-6+1=0$ $\checkmark$), donc $H=M$.
\end{exemple}

\begin{exemple}[Exercice 4 -- coplanarité par produit mixte]
Montrer que les points $A(1,1,0)$, $B(2,0,1)$, $C(0,1,1)$, $D(1,0,2)$ sont coplanaires.

\smallskip
\textbf{Corrigé.} $\overrightarrow{AB}(1,-1,1)$, $\overrightarrow{AC}(-1,0,1)$, $\overrightarrow{AD}(0,-1,2)$.
\[
  \overrightarrow{AB}\wedge\overrightarrow{AC} = ((-1)(1)-(1)(0)\,;\ (1)(-1)-(1)(1)\,;\ (1)(0)-(-1)(-1)) = (-1,-2,-1).
\]
\[
  (\overrightarrow{AB}\wedge\overrightarrow{AC})\cdot\overrightarrow{AD} = (-1)(0)+(-2)(-1)+(-1)(2) = 0+2-2 = 0.
\]
Le produit mixte est nul : les quatre points sont \textbf{coplanaires}.
\end{exemple}

% ------------------- SECTION 5 -------------------
\section{Méthodes et pièges : la boîte à outils de l'examen}

\subsection{Ce qu'on te demande, ce que tu fais}

\begin{center}
\renewcommand{\arraystretch}{1.45}
\sffamily\small
\begin{tabular}{>{\raggedright\arraybackslash}p{7.2cm} >{\raggedright\arraybackslash}p{8cm}}
\rowcolor{qtealdark} \color{white}\bfseries Ce qu'on te demande & \color{white}\bfseries Ton réflexe \\
\rowcolor{tealpale} Montrer $(AB)\perp(AC)$ & $\overrightarrow{AB}\cdot\overrightarrow{AC}=0$. \\
Trouver l'équation d'un plan (point + normal) & $a(x-x_0)+b(y-y_0)+c(z-z_0)=0$. \\
\rowcolor{tealpale} Trouver le plan $(ABC)$ (3 points) & Vecteur normal $=\overrightarrow{AB}\wedge\overrightarrow{AC}$, puis équation avec $A$. \\
Distance point-plan & $d=\dfrac{|ax_0+by_0+cz_0+d|}{\sqrt{a^2+b^2+c^2}}$. \\
\rowcolor{tealpale} Aire d'un triangle dans l'espace & $\mathcal{A}=\frac12\|\overrightarrow{AB}\wedge\overrightarrow{AC}\|$. \\
Volume d'un tétraèdre & $V=\frac16|(\overrightarrow{AB}\wedge\overrightarrow{AC})\cdot\overrightarrow{AD}|$. \\
\rowcolor{tealpale} Montrer 4 points coplanaires & Produit mixte nul. \\
Tester la colinéarité de deux vecteurs & $\vec u\wedge\vec v=\vec0$ (ou déterminant $2\times2$ dans le plan). \\
\end{tabular}
\end{center}

\subsection{Les pièges qui coûtent des points}

\begin{remarque}[Erreurs relevées chaque année par les correcteurs]
\begin{itemize}[leftmargin=6mm, itemsep=0.5mm]
  \item Confondre produit scalaire (donne un \textbf{réel}) et produit vectoriel (donne un \textbf{vecteur}).
  \item Erreur de signe dans les coordonnées du produit vectoriel : bien retenir l'ordre $yz'-zy'\,;\,zx'-xz'\,;\,xy'-yx'$ (permutation circulaire).
  \item Oublier la valeur absolue dans la formule de distance point-plan.
  \item Oublier le facteur $\frac12$ (aire du triangle) ou $\frac16$ (volume du tétraèdre) par rapport au parallélogramme/parallélépipède.
  \item Utiliser le produit vectoriel dans le \textbf{plan} : il n'y est pas défini de la même façon (c'est une notion propre à l'espace, ou au déterminant $2\times2$ pour la colinéarité plane).
  \item Confondre vecteur \textbf{directeur} d'une droite et vecteur \textbf{normal} d'un plan.
\end{itemize}
\end{remarque}

% ------------------- RESUME -------------------
\vspace{4mm}
\begin{tcolorbox}[
  enhanced, boxrule=0pt, arc=2mm,
  interior style={left color=qtealdark, right color=qteal},
  left=6mm, right=6mm, top=4mm, bottom=4mm,
  fontupper=\color{white}\sffamily,
  title={\sffamily\bfseries\color{white}\ding{72}~~À retenir},
  colbacktitle=qtealdark, coltitle=white, boxrule=0pt
]
\begin{itemize}[leftmargin=6mm, label={\color{qamber}\ding{212}}, itemsep=1mm]
  \item $\vec u\cdot\vec v=xx'+yy'+zz'$ (repère orthonormé) ; $\vec u\perp\vec v\iff\vec u\cdot\vec v=0$.
  \item Plan de normal $(a,b,c)$ passant par $A$ : $a(x-x_0)+b(y-y_0)+c(z-z_0)=0$ ; distance $d=\frac{|ax_0+by_0+cz_0+d|}{\sqrt{a^2+b^2+c^2}}$.
  \item $\vec u\wedge\vec v=(yz'-zy'\,;zx'-xz'\,;xy'-yx')$ ; orthogonal à $\vec u$ et $\vec v$ ; nul ssi colinéaires.
  \item Aire triangle $=\frac12\|\overrightarrow{AB}\wedge\overrightarrow{AC}\|$ ; volume tétraèdre $=\frac16|(\overrightarrow{AB}\wedge\overrightarrow{AC})\cdot\overrightarrow{AD}|$ ; coplanarité $\iff$ produit mixte nul.
\end{itemize}
\end{tcolorbox}

\end{document}
